% Conteudo migrado de thesis/main.tex.
% O titulo do capitulo fica definido em thesis.tex.

\section{Microcontrolador}

A primeira decisão de hardware foi a escolha do microcontrolador para o nó sensor.
A opção mais natural inicialmente seria utilizar a família STM32, já presente
no sistema de aquisição da equipe e fornecida por um patrocinador. % TODO: confirmar patrocínio

A linha STM32 não possui nenhum chip com Wi-Fi nativo. A única família com comunicação
sem fio integrada é a STM32WB, que oferece Bluetooth Low Energy e 802.15.4 —
sem suporte a Wi-Fi. O STM32WB integra um núcleo Cortex-M4 para a aplicação
e um Cortex-M0+ dedicado ao stack de rádio, sendo uma plataforma robusta, bem
documentada e familiar à equipe \cite{st_stm32wb_series} — uma opção forte, desde que o BLE atendesse
aos requisitos de comunicação.

Para contornar a ausência de Wi-Fi nativo, uma alternativa seria combinar um
STM32 convencional com um módulo Wi-Fi externo. No entanto, essa abordagem exigiria
três componentes distintos na placa do nó sensor: o microcontrolador, o módulo
de comunicação e um circuito de gerenciamento de bateria separado. A
combinação resulta em uma solução com footprint físico, complexidade de integração
e custo de desenvolvimento incompatíveis com os requisitos de compacidade e leveza
do projeto, e foi descartada por esse motivo.

A terceira opção considerada foi o ESP32, plataforma da Espressif que integra
nativamente Wi-Fi e Bluetooth em um único chip. Módulos de desenvolvimento
compactos e com gerenciamento de bateria integrado estão amplamente
disponíveis para essa plataforma — como o Beetle ESP32-C3 da DFRobot, de
25×20,5 mm, que inclui gerenciamento de bateria na própria placa, sem
necessidade de componentes externos adicionais \cite{dfrobot_beetle_esp32c3}.

A decisão entre as duas plataformas foi tomada em conjunto com a escolha do protocolo
de comunicação, discutida a seguir.

% === SECTION: Protocolo de Comunicação Sem Fio ===

\section{Protocolo de Comunicação Sem Fio}

As opções naturais de protocolo para a solução sem fio foram Bluetooth Low Energy, Wi-Fi convencional e, posteriormente, ESP-NOW.

O BLE, disponível tanto no STM32WB quanto no ESP32, foi a primeira opção considerada. No entanto, o protocolo apresenta limitações importantes frente aos requisitos do sistema. O estabelecimento de uma conexão BLE envolve descoberta, pareamento ou associação entre dispositivos, o que introduz latência inicial e estado de conexão incompatíveis com a ideia de um sensor que pode começar a transmitir imediatamente. Além disso, o gerenciamento de múltiplas conexões simultâneas adicionaria complexidade significativa à aplicação. Essas limitações descartaram o BLE como protocolo principal e, por consequência, reduziram a atratividade do STM32WB para este projeto. Comparações experimentais entre ESP-NOW, Wi-Fi e Bluetooth também indicam que diferenças de latência, alcance e consumo tornam a escolha do protocolo dependente do perfil de aplicação \cite{eridani2021_espnow_wifi_bluetooth}.

O Wi-Fi convencional resolveria parte do problema de alcance e vazão, mas manteria o overhead de associação a uma rede, configuração de infraestrutura, obtenção de endereço IP e estabelecimento de sessão. Esse comportamento é inadequado para um sistema que deve funcionar sem ponto de acesso, sem configuração manual de rede e com início de transmissão o mais imediato possível.

Durante o desenvolvimento, foi adotado o ESP-NOW, protocolo proprietário da Espressif para comunicação direta entre dispositivos ESP32. O ESP-NOW opera sobre a camada física do Wi-Fi, mas não exige associação a uma rede Wi-Fi convencional. Os quadros são endereçados diretamente por endereço MAC, utilizam a verificação de integridade do próprio quadro 802.11, incluindo CRC em nível de enlace, e, na versão clássica utilizada neste trabalho, permitem até 250 bytes de payload por transmissão \cite{espressif_espnow_programming_guide,pasic2021_espnow_esp32}. Essa combinação atende aos requisitos de baixo overhead, operação sem infraestrutura e comunicação entre múltiplos dispositivos.

A principal limitação introduzida pelo ESP-NOW aparece no uso de broadcast. Quando um pacote é enviado para o endereço \texttt{FF:FF:FF:FF:FF:FF}, o transmissor não recebe uma confirmação nativa de entrega associada a cada receptor. Portanto, qualquer garantia de confiabilidade precisa ser construída acima do ESP-NOW, na camada de aplicação. Essa restrição motivou o desenvolvimento de duas variantes internas do transporte WCAN sobre o mesmo protocolo físico: uma variante broadcast, que preserva a semântica de barramento compartilhado do CAN, e uma variante unicast, que usa descoberta de receptores e confirmação do driver Wi-Fi por destinatário. As duas variantes implementam a mesma interface de aplicação e foram mantidas em paralelo para permitir comparação experimental em condições realistas, sem escolher antecipadamente uma solução definitiva.

Esse trade-off foi considerado aceitável porque o objetivo do projeto é expandir a aquisição de dados do veículo, não substituir por rádio uma malha de controle crítica. A camada WCAN pode acrescentar acknowledgments, retransmissões, filtragem, descoberta de receptores e medição de perdas, enquanto o ESP-NOW fornece a transmissão imediata e sem infraestrutura que os requisitos exigem. A escolha entre broadcast e unicast, portanto, deve ser guiada pelos resultados de entrega, latência e robustez observados nos ensaios, como também ocorre em outras soluções de telemetria veicular baseadas em Wi-Fi ou ESP32 \cite{fernandes2022_formula_sae_wifi_can}.
